\documentclass[11pt,a4paper]{article}
\usepackage[utf8]{inputenc}
\usepackage[T1]{fontenc}
\usepackage[french]{babel}
\usepackage{amsmath,amssymb}
\usepackage[margin=2.4cm]{geometry}
\usepackage{booktabs}
\usepackage{hyperref}
\hypersetup{colorlinks=true,linkcolor=blue,citecolor=blue,urlcolor=blue}
\usepackage{natbib}

\title{\bf L'exposant de dilution de la matière n'est pas\\
identifiable à partir de priors CMB comprimés :\\
une ambiguïté de référent, pas une limite de précision}
\author{Édouard Lantenois\thanks{\texttt{github.com/Dantenos}}}
\date{Août 2026}

\begin{document}
\maketitle

\begin{abstract}
\noindent
Considérons la famille à un paramètre où la matière noire dilue en
$\rho_{\rm dm}\propto a^{-3+\epsilon}$ avec une véritable constante cosmologique.
L'identité de Bianchi impose que ce fond soit celui d'un fluide d'équation d'état constante
$w=-\epsilon/3$ : la famille n'est donc pas une alternative au $\Lambda w$DM mais sa
réécriture. Or les priors de distance CMB comprimés $(\mathcal{R}, \ell_A, \omega_b)$, ainsi
que les formules d'ajustement usuelles pour $r_d$, $r_\ast$ et $z_\ast$, réclament tous
\emph{une seule} valeur de $\omega_m$. Quand $\epsilon\neq0$ il y en a deux --- l'étiquette
d'aujourd'hui et la valeur à la recombinaison --- et rien dans les données ne désigne
laquelle. Nous montrons, sur des données identiques (Pantheon+, BAO DESI DR2, priors de
distance Planck 2018, SH0ES), que ce seul choix de comptabilité déplace l'exposant recouvré
de $\epsilon=+0{,}0070$ à $\epsilon=-0{,}0120$, soit une étendue de $0{,}019$, quand chaque
configuration annonce $\sigma\simeq0{,}002$. Nous montrons ensuite que ce n'est
\emph{pas} un effet de perte d'information : les priors comprimés contraignent $\omega_c$ à
$\sigma=0{,}0010$, valeur identique à celle que rend la vraisemblance Planck complète
TT,TE,EE$+$bas-$\ell$ à paramètres fixés identiques --- un rapport de $1{,}00$, borné à
$]0{,}5\,;1{,}25[$ par la quantification de grille. L'ambiguïté est donc de \emph{référent},
non de \emph{précision}, et l'étendue vaut dix-neuf fois la précision disponible. De façon
cohérente, les déterminations publiées de cet exposant obtenues avec un code de Boltzmann
modifié et une vraisemblance CMB complète se groupent à $\epsilon = -0{,}00064\pm0{,}00045$
et s'étendent sur $0{,}005$, quand celles par priors comprimés s'étendent sur $0{,}016$. Nous
rapportons un piège supplémentaire : appliquer $\epsilon$ à la matière \emph{totale} plutôt
qu'à la seule matière noire crée une incohérence interne valant $7{,}5$ à $15{,}1\sigma$ sur
$\ell_A$, la même vraisemblance traitant alors les baryons comme standard dans $R_b$ et dans
le prior sur $\omega_b$. Nous ne revendiquons aucune détection. Nous établissons qu'un
$\epsilon$ publié à partir de priors comprimés porte une systématique non citée de l'ordre de
$10^{-2}$ tant que le choix de $\omega_m$ n'est pas énoncé.
\end{abstract}

\section{La famille, et pourquoi elle n'a qu'un fond}
\label{sec:famille}

Écrivons le secteur sombre comme une composante de matière dont la densité obéit à
\begin{equation}
\rho_{\rm dm}(a) = \rho_{{\rm dm},0}\,a^{-3+\epsilon},
\qquad
\rho_\Lambda = \text{const},
\qquad
\rho_b \propto a^{-3},\quad \rho_r\propto a^{-4}.
\label{eq:famille}
\end{equation}
Il est tentant de lire \eqref{eq:famille} comme de la \emph{matière sans pression non
conservée}, le déficit étant absorbé par rien en particulier. Cette lecture n'est pas
disponible. Avec $\Lambda$ véritablement constante, $\nabla_\mu T^{\mu\nu}=0$ appliqué au
tenseur total \emph{impose} que le secteur de matière porte une pression
\begin{equation}
p_{\rm dm} = -\tfrac{\epsilon}{3}\,\rho_{\rm dm},
\qquad\text{c'est-à-dire}\qquad
w_{\rm dm} = -\epsilon/3 .
\label{eq:bianchi}
\end{equation}
« Matière sans pression non conservée avec $\Lambda$ constante » et « matière conservée
d'équation d'état constante $w=-\epsilon/3$ avec $\Lambda$ constante » sont le même modèle
écrit deux fois. Au niveau du fond il y a un modèle, pas deux.

Ce n'est pas neuf. \citet{lima1996} donnent, pour la création adiabatique de matière de
paramètre $\beta$, l'indice adiabatique effectif $\gamma_\ast=\gamma(1-\beta)$ et la loi
$\rho=\rho_0(a_0/a)^{3\gamma(1-\beta)}$, avec une pression de création $p_c=-\beta\gamma\rho$.
Pour la poussière ($\gamma=1$) c'est exactement
\eqref{eq:famille}--\eqref{eq:bianchi} avec $\epsilon=3\beta$. Deux mises en garde, que nous
énonçons parce qu'elles se perdent facilement : l'identité $\gamma_\ast=\gamma(1-\beta)$ ne
vaut que pour $\beta$ \emph{constant}, et la correspondance $\epsilon=3\beta$ est
\emph{dépendante de l'espèce} --- pour le rayonnement ($\gamma=4/3$) le même $\beta$ donne
$\rho_r\propto a^{-4(1-\beta)}$. On ne peut pas la transporter à travers l'ère de rayonnement.

Notons aussi que le symbole $\epsilon$ est employé dans une littérature voisine mais
différente --- celle du \emph{vide décroissant}, par exemple \citet{alcaniz2005} --- où
$\Lambda$ n'est pas constante. La famille \eqref{eq:famille} est celle à $\Lambda$ constante,
et son fond coïncide avec le modèle $\Lambda w$DM contraint par \citet{kumar2025}.

\section{Où l'ambiguïté entre}
\label{sec:ou}

Les analyses comprimées de cette famille emploient trois ingrédients, et chacun d'eux prend
$\omega_m$ en entrée :
\begin{enumerate}
\item le paramètre de décalage $\mathcal{R}=\sqrt{\Omega_m H_0^2}\,D_A(z_\ast)/c$, dont le
      préfacteur $\sqrt{\Omega_m}$ est dérivé en supposant $\rho_m\propto a^{-3}$ ;
\item les formules d'ajustement des horizons sonores $r_d$, $r_\ast$ \citep{eh1998},
      calibrées sur des grilles $\Lambda$CDM ;
\item la formule d'ajustement de $z_\ast$ \citep{hs1996}.
\end{enumerate}
En $\Lambda$CDM l'entrée est sans ambiguïté. Dans \eqref{eq:famille} elle ne l'est pas :
l'étiquette $\Omega_m$ et la densité qui gouverne réellement la physique d'avant
recombinaison diffèrent de
\begin{equation}
\frac{\Omega_m^{\rm rec}}{\Omega_m^{\rm étiq}} = a_\ast^{\,\epsilon},
\qquad a_\ast\simeq 1/1091,
\qquad \ln a_\ast \simeq -6{,}995 .
\label{eq:ratio}
\end{equation}
Pour $|\epsilon|=0{,}005$ l'écart vaut $3{,}6\%$. L'analyste doit choisir, et la littérature
ne dit pas lequel. \citet{elgaroy2007} avertissaient que $\mathcal{R}$ « n'est pas
directement mesurable à partir du fond diffus, et que sa valeur est habituellement dérivée
des données en supposant une cosmologie plate avec matière noire et constante
cosmologique » ; nous chiffrons ce que coûte cet avertissement dans cette famille.

\section{La mesure : quatre choix de comptabilité, un seul jeu de données}
\label{sec:mesure}

Nous employons Pantheon+ (1580 SNe, $z>0{,}01$, calibrateurs exclus), les 13 mesures BAO de
DESI DR2, les priors de distance Planck 2018 et le prior SH0ES sur $H_0$. Quatre
configurations ne diffèrent que par la comptabilité :

\begin{itemize}
\item[(a)] \textbf{étiquette} : les formules d'ajustement et $\mathcal{R}$ reçoivent le
  $\Omega_m$ d'aujourd'hui ;
\item[(b)] \textbf{cohérente} : elles reçoivent $\Omega_m a_\ast^{\epsilon}$, la densité à la
  recombinaison ;
\item[(c)] \textbf{directe, sans $\mathcal{R}$} : $r_s$ est obtenu en intégrant
  $\int c_s\,\mathrm{d}a/(a^2H)$ dans le fond propre du modèle, et $\mathcal{R}$ est retiré,
  laissant $(\ell_A,\omega_b)$ avec la covariance marginale $2\times2$ ;
\item[(d)] \textbf{directe, $\mathcal{R}$ cohérent} : comme (c) mais en conservant
  $\mathcal{R}$ alimenté par la densité de recombinaison.
\end{itemize}

L'intégrateur direct est validé contre la table CAMB qu'il remplace : avec
$N_{\rm eff}=3{,}046$ et le retrait de la contribution des neutrinos massifs du budget de
matière, le rapport $r_s^{\rm direct}/r_s^{\rm CAMB}$ varie de $1{,}3\times10^{-6}$ sur toute
la grille $(\omega_b,\omega_m)$, soit $0{,}01\sigma$ sur $\ell_A$. Le décalage de traînée
qu'il infère, récupéré en inversant l'intégrateur contre le $r_{\rm drag}$ de CAMB, tombe dans
$[1057{,}8\,;1062{,}1]$ contre $1059{,}9$ chez Planck --- un contrôle que nous n'avions pas
imposé.

\begin{table}[t]
\centering
\begin{tabular}{lrrr}
\toprule
configuration & $\epsilon$ & $\sigma$ (opposable) & $\Delta\chi^2$ contre $\Lambda$CDM \\
\midrule
(a) étiquette                 & $+0{,}0070$ & $2{,}33$ & $+9{,}36$ \\
(b) cohérente                 & $-0{,}0040$ & $2{,}00$ & $+9{,}36$ \\
(c) directe, sans $\mathcal{R}$ & $-0{,}0120$ & $3{,}14$ & $+9{,}87$ \\
(d) directe, $\mathcal{R}$ cohérent & $-0{,}0040$ & $2{,}00$ & $+9{,}90$ \\
\midrule
\textbf{étendue} & \textbf{0{,}0190} & & \\
\bottomrule
\end{tabular}
\caption{Mêmes données, même modèle, quatre choix de comptabilité. Le $\sigma$ cité est le
plus petit de la demi-largeur du profil et de la valeur de Wilks $\sqrt{\Delta\chi^2}$ ;
lorsque les deux divergent, le profil n'est pas parabolique et le $\sigma$ local n'est pas
opposable.}
\label{tab:quatre}
\end{table}

La table~\ref{tab:quatre} est la mesure centrale. L'étendue vaut $0{,}019$ ; chaque
configuration annonce $\sigma\simeq0{,}002$. Deux remarques d'honnêteté. D'abord, la
configuration (c) rapporte une significativité locale de $6{,}0\sigma$ que nous \emph{ne}
citons \emph{pas} : son profil est aplati ($\Delta\chi^2(2\sigma_{\rm local})/4 = 0{,}39$ au
lieu de $1$) et la valeur par rapport de vraisemblance vaut $3{,}14$. Ensuite, la
décomposition de son gain montre que $126\%$ de celui-ci vient du seul terme SH0ES, les
supernovae empirant de $1{,}72$ et les BAO de $0{,}96$ : cette configuration achète $H_0$,
elle ne mesure pas la dilution. Nous la rapportons parce qu'elle appartient à l'étendue, non
parce que nous y croyons.

\section{Ce n'est pas une perte d'information}
\label{sec:pouvoir}

L'explication naturelle --- les priors comprimés ne portent que de la géométrie, la
vraisemblance complète porte aussi la \emph{hauteur} des pics acoustiques qui épingle
$\omega_m$ --- est fausse, et nous l'avons testée plutôt qu'affirmée.

La prémisse est juste : évaluer le $\chi^2$ comprimé à $\ln(10^{10}A_s)=3{,}000$ et
$3{,}100$ rend $2{,}575399113$ les deux fois, à la précision machine. Il ne porte
littéralement aucune information d'amplitude.

La conclusion ne suit pas. En profilant $\omega_c$ avec CAMB et la vraisemblance Planck
TT,TE,EE$+$bas-$\ell$ à $\theta_\ast$, $\omega_b$, $n_s$, $\tau$ fixés, et en profilant
$\ln(10^{10}A_s)$ en chaque point, on obtient $\sigma(\omega_c)=0{,}0010$. Le même profil à
travers les priors comprimés donne $\sigma(\omega_c)=0{,}0010$. Le rapport vaut $1{,}00$ ;
en propageant la quantification de grille il est borné dans $]0{,}5\,;1{,}25[$, très loin de
toute signature de perte d'information. En contrôle de cohérence, notre $\sigma$ est $1{,}20$
fois plus petit que la valeur marginalisée $\omega_c=0{,}1200\pm0{,}0012$ de
\citet{planck2018}, ce qu'on attend en figeant quatre paramètres ; et le $\chi^2$ complet à
son optimum $\Lambda$CDM reproduit notre ancre gelée à $0{,}005$ près.

Les deux déterminations de $\omega_m$ sont donc également fines --- $0{,}7\%$ --- et
l'étendue entre choix de comptabilité vaut $0{,}019$, soit \emph{dix-neuf fois} la précision
disponible.

\begin{equation}
\boxed{\;\text{L'ambiguïté est de référent, non de précision.}\;}
\end{equation}

$\mathcal{R}$ et les formules d'ajustement réclament un $\omega_m$ ; la famille en fournit
deux ; les données n'en départagent aucun. Quatre analystes disposant des mêmes données et de
contraintes identiques à $0{,}7\%$ obtiennent des réponses étalées sur $0{,}019$ --- non par
bruit, mais par choix.

\section{Le piège baryonique}
\label{sec:baryons}

Un second mode de défaillance mérite d'être isolé, parce qu'on y tombe facilement et qu'il est
grand. Si $\epsilon$ est appliqué à la matière \emph{totale} --- c'est-à-dire si $\Omega_m$
dans \eqref{eq:famille} se lit baryons plus matière noire --- alors les baryons diluent aussi
anormalement. Mais la même vraisemblance, typiquement, (i) compare $\omega_b$ à un prior
mesuré sous $\rho_b\propto a^{-3}$ et (ii) intègre l'horizon sonore avec
$R_b=(3\omega_b/4\omega_\gamma)\,a$, forme qui suppose la même chose. La cohérence exigerait
$R_b\propto a^{1+\epsilon}$. Mesurons l'écart :

\begin{table}[h]
\centering
\begin{tabular}{lrr}
\toprule
 & $|\epsilon|=0{,}003$ & $|\epsilon|=0{,}006$ \\
\midrule
$r_s$ sous $R_b\propto a$ contre $a^{1+\epsilon}$ & $7{,}5\sigma$ sur $\ell_A$ & $15{,}1\sigma$ \\
$\omega_b$(recombinaison) contre le prior Planck & $4{,}8\sigma$ & $8{,}1\sigma$ \\
\bottomrule
\end{tabular}
\caption{Coût d'appliquer $\epsilon$ à la matière totale tout en traitant les baryons comme
standard ailleurs dans la même vraisemblance.}
\label{tab:baryons}
\end{table}

Les deux dépassent tout ce que l'analyse cherche à mesurer. Le remède est de restreindre
$\epsilon$ au secteur sombre, comme le font \citet{kumar2025} et \citet{yaoliu2025}. Nous
notons, contre notre propre confort, que l'argument souvent invoqué « la nucléosynthèse
primordiale l'interdit » est plus faible qu'il n'y paraît : exiger que les déterminations BBN
et CMB de $\omega_b$ s'accordent à $\sim2{,}5\%$ sur les $\simeq12{,}5$ e-plis séparant les
deux époques ne borne que $|\epsilon_b|\lesssim2\times10^{-3}$, et nous n'avons trouvé aucune
contrainte publiée sur l'exposant de dilution des baryons lui-même. Le défaut ci-dessus est
une incohérence \emph{interne}, pas la violation d'une borne extérieure.

\section{Confrontation aux déterminations publiées}
\label{sec:litterature}

Nous avons rassemblé toute détermination de cet exposant vérifiable à la source primaire, en
convertissant chacune vers la convention de \eqref{eq:famille} par une règle déclarée, et en
les séparant par méthode.

\begin{table}[t]
\centering
\begin{tabular}{llrr}
\toprule
détermination & méthode & $\epsilon$ & $\sigma$ \\
\midrule
\multicolumn{4}{l}{\emph{vraisemblance complète, code de Boltzmann modifié}}\\
\citet{tsiapi2019}, $\Lambda(H)$CDM2, Planck seul & CAMB, Planck 2015 & $+0{,}00148$ & $0{,}00252$\\
\citet{tsiapi2019}, $\Lambda(H)$CDM2, joint       & CAMB, Planck 2015 & $-0{,}00020$ & $0{,}00181$\\
\citet{tsiapi2019}, $\Lambda(H)$CDM1, joint$^\dagger$ & CAMB, Planck 2015 & $+0{,}00302$ & $0{,}00151$\\
\citet{kumar2025}, PL18$+$DESI DR2                & CAMB, Planck 2018 & $-0{,}00194$ & $0{,}00096$\\
\citet{li2025}, Planck$+$DESI DR2$+$DESY5         & CAMB, NPIPE       & $-0{,}00126$ & $0{,}00075$\\
\citet{li2024}, CMB$+$DESI DR1$+$DESY5$^\ddagger$ & CosmoMC           & $-0{,}00025$ & $0{,}00092$\\
\midrule
\multicolumn{2}{l}{\textbf{moyenne pondérée}} & $\mathbf{-0{,}00064}$ & $\mathbf{0{,}00045}$\\
\multicolumn{2}{l}{étendue} & \multicolumn{2}{l}{$0{,}0050$}\\
\midrule
\multicolumn{4}{l}{\emph{priors de distance comprimés}}\\
\citet{yang2025}, cinq jeux de données            & priors, DESI DR1 & $-0{,}00612$ & $0{,}00243$\\
ce travail, configuration (a)                     & priors           & $+0{,}0060$ & $0{,}0030$\\
ce travail, configuration (b)                     & priors           & $-0{,}0030$ & $0{,}0010$\\
ce travail, configuration (c)                     & priors           & $-0{,}0100$ & $0{,}0020$\\
ce travail, $r_s$ direct, $\mathcal{R}$ étiquette & priors           & $+0{,}0020$ & $0{,}0020$\\
\midrule
\multicolumn{2}{l}{étendue} & \multicolumn{2}{l}{$\mathbf{0{,}0160}$}\\
\bottomrule
\end{tabular}
\caption{$^\dagger$ Conversion approchée : le $Q$ de ce modèle est aussi alimenté par les
baryons et le rayonnement. Nous l'incluons parce qu'il \emph{élargit} la famille en
vraisemblance complète et rend donc notre propre démonstration plus difficile.
$^\ddagger$ Prédécesseur emboîté de \citet{li2025}, même groupe ; inclus et signalé.
Conversions : $\epsilon=3\nu f$ ou $\epsilon=-3w_{\rm dm}f$ avec
$f=\rho_{\rm dm}/\rho_m=0{,}8389$.}
\label{tab:litterature}
\end{table}

La famille en vraisemblance complète est cohérente en interne
($\chi^2/\mathrm{ddl}=1{,}87$ pour cinq degrés de liberté) et s'étend sur $0{,}0050$ ; celle
par priors comprimés s'étend sur $0{,}0160$, un facteur $3{,}23$. En excluant la conversion
approchée signalée --- le choix confortable, que nous ne prenons pas --- le facteur vaudrait
$4{,}67$ et $\chi^2/\mathrm{ddl}=0{,}72$.

Trois réserves accordées entièrement. Les six modèles en vraisemblance complète ne sont pas
identiques à \eqref{eq:famille} : \citet{tsiapi2019} font courir $\rho_\Lambda$ et portent un
préfacteur $\Omega_{\rm dm}/(1-\nu)$ dans $E^2$, de sorte qu'à $\epsilon$ égal leur $H(z)$
n'est pas le nôtre ; \citet{kumar2025} et \citet{yaoliu2025} attachent des prescriptions de
perturbations différentes ($c_s^2=0$ contre $c_s^2=w$) au même fond, et leurs intervalles ne
sont pas transférables à une loi de fond seul. Leurs données diffèrent des nôtres et entre
elles. Et deux des six sont du même groupe sur des données emboîtées. La
table~\ref{tab:litterature} mesure donc une \emph{dispersion}, et c'est tout ce que nous en
tirons.

Nous consignons aussi, parce que c'est l'élément le plus instructif de la colonne comprimée :
le résultat phare de \citet{yang2025}, $\epsilon=-0{,}0073^{+0{,}0029}_{-0{,}0033}$ à
$2{,}4\sigma$, exige cinq jeux de données ; sur DESI$+$CMB$+$CC seuls le même article rapporte
$\epsilon=+0{,}0023^{+0{,}0055}_{-0{,}0067}$ et conclut à l'absence de déviation
significative. Un renversement de signe à l'intérieur d'une seule publication, sur des données
emboîtées.

\section{Ce que nous ne revendiquons pas}
\label{sec:limites}

Nous ne mesurons pas $\epsilon$. Aucune de nos quatre configurations n'est une vraisemblance
complète, et celle qui retire $\mathcal{R}$ est pilotée par le prior sur $H_0$. Nous ne
réfutons ni \citet{yang2025} ni personne : notre pipeline, nos données et notre traitement des
nuisances diffèrent. Nous ne prétendons pas que le domaine ait abandonné $\mathcal{R}$ ---
DESI DR2 conditionne sa ligne de base sur $(\theta_\ast,\omega_b,\omega_{bc})$ mais teste
explicitement la compression $(\mathcal{R},\ell_A,\omega_b)$ et la déclare équivalente pour
les modèles tardifs qu'ils considèrent \citep{desidr2}. Retirer $\mathcal{R}$ se justifie ici
étroitement, parce que $\sqrt{\Omega_m}$ n'a pas de définition quand la matière ne dilue pas
en $a^{-3}$.

Nous divulguons également une systématique qui nous est propre. Les priors de distance que
nous employons sont ceux de \citet{zhaiwang2019}, dont le $\ell_A=301{,}80845$ correspond à
$100\theta=1{,}040923$, c'est-à-dire la convention $\theta_{\rm MC}$. Notre $\chi^2$ prédit
$\ell_A$ à partir des $z_\ast$ et $r_\ast$ exacts de CAMB, donc en convention $\theta_\ast$,
rendant $\ell_A=301{,}75963$ et $100\theta=1{,}041091$ à l'optimum $\Lambda$CDM. Les deux
conventions diffèrent de $-0{,}54\sigma$ sur $\ell_A$ à chaque évaluation. Ce décalage est
commun aux quatre configurations et n'engendre donc pas l'étendue de la
table~\ref{tab:quatre}, mais il est présent et nous l'énonçons plutôt que de laisser un
lecteur le découvrir.

\section{Recommandations}
\label{sec:reco}

\begin{enumerate}
\item Une contrainte publiée sur un exposant de dilution de la matière obtenue par priors CMB
      comprimés devrait énoncer \emph{quel} $\omega_m$ alimente $\mathcal{R}$, $r_d$,
      $r_\ast$ et $z_\ast$, et citer la systématique obtenue en faisant l'autre choix. Dans
      cette famille elle est de l'ordre de $10^{-2}$ --- un ordre de grandeur au-dessus de
      l'erreur statistique habituellement citée.
\item $\epsilon$ devrait n'être appliqué qu'au secteur sombre, et l'intégrande de l'horizon
      sonore devrait employer la même loi baryonique que le fond (\S\ref{sec:baryons}).
\item Là où la voie comprimée est conservée, $r_s$ devrait être intégré dans le fond propre du
      modèle plutôt que tiré d'une formule d'ajustement calibrée sur $\Lambda$CDM. C'est la
      pratique courante dans la littérature de l'énergie sombre précoce et cela ne coûte
      rien ; nous validons une implémentation à $1{,}3\times10^{-6}$ contre CAMB au
      \S\ref{sec:mesure}.
\item Enfin, nous consignons une lacune. Un balayage systématique de la littérature de la
      création de matière n'a rendu \emph{aucune} analyse d'un tel modèle contre les spectres
      de puissance CMB complets : toutes emploient des supernovae, des BAO, $H(z)$, des
      données de croissance, ou un CMB comprimé. L'exemple le plus récent que nous ayons
      trouvé décrit encore son entrée CMB comme « une vraisemblance CMB comprimée (dite aussi
      géométrique ou de fond), où le CMB est traité comme une mesure BAO à $z\approx1100$ »
      \citep{schiavone2026}. Au vu du \S\ref{sec:pouvoir}, le coût de ce choix n'est pas une
      perte de précision --- c'est l'introduction d'un référent que les données ne peuvent pas
      fixer.
\end{enumerate}

\section*{Données et reproductibilité}
Chaque nombre de cet article est produit par un script dont les critères de succès ont été
gelés par empreinte avant exécution, puis revérifiés après coup par un outil indépendant
lui-même testé par mutation (cinq fautes injectées, cinq détectées). Là où un critère gelé
s'est révélé mal posé, le critère a été laissé intact et le défaut consigné, plutôt que le
critère réécrit.

\begin{thebibliography}{99}
\bibitem[Alcaniz \& Lima(2005)]{alcaniz2005} Alcaniz, J.~S., \& Lima, J.~A.~S. 2005,
  Phys.\ Rev.\ D, 72, 063516, \texttt{astro-ph/0507372}
\bibitem[DESI Collaboration(2025)]{desidr2} DESI Collaboration 2025, Phys.\ Rev.\ D, 112,
  083515, \texttt{arXiv:2503.14738}
\bibitem[Eisenstein \& Hu(1998)]{eh1998} Eisenstein, D.~J., \& Hu, W. 1998, ApJ, 496, 605
\bibitem[Elgar\o y \& Multam\"aki(2007)]{elgaroy2007} Elgar\o y, \O., \& Multam\"aki, T.
  2007, A\&A, 471, 65, \texttt{astro-ph/0702343}
\bibitem[Hu \& Sugiyama(1996)]{hs1996} Hu, W., \& Sugiyama, N. 1996, ApJ, 471, 542
\bibitem[Kumar et al.(2025)]{kumar2025} Kumar, U., Ajith, A., \& Verma, A. 2025,
  \texttt{arXiv:2504.14419}
\bibitem[Li et al.(2024)]{li2024} Li, T.-N., Wu, P.-J., Du, G.-H., et al. 2024, ApJ, 976, 1,
  \texttt{arXiv:2407.14934}
\bibitem[Li et al.(2025)]{li2025} Li, T.-N., et al. 2025, \texttt{arXiv:2510.11363}
\bibitem[Lima et al.(1996)]{lima1996} Lima, J.~A.~S., Germano, A.~S.~M., \& Abramo, L.~R.~W.
  1996, Phys.\ Rev.\ D, 53, 4287, \texttt{gr-qc/9511006}
\bibitem[Planck Collaboration(2020)]{planck2018} Planck Collaboration 2020, A\&A, 641, A6,
  \texttt{arXiv:1807.06209}
\bibitem[Schiavone et al.(2026)]{schiavone2026} Schiavone, T., De Angelis, M., Escamilla,
  L.~A., Montani, G., \& Di Valentino, E. 2026, \texttt{arXiv:2601.14222}
\bibitem[Tsiapi \& Basilakos(2019)]{tsiapi2019} Tsiapi, P., \& Basilakos, S. 2019, MNRAS,
  485, 2505, \texttt{arXiv:1810.12902}
\bibitem[Yang et al.(2025)]{yang2025} Yang, W., Dai, W.-B., \& Wang, D. 2025,
  \texttt{arXiv:2505.09879}
\bibitem[Yao \& Liu(2025)]{yaoliu2025} Yao, Y., \& Liu, X. 2025, \texttt{arXiv:2507.00478}
\bibitem[Zhai \& Wang(2019)]{zhaiwang2019} Zhai, Z., \& Wang, Y. 2019, JCAP, 07, 005,
  \texttt{arXiv:1811.07425}
\end{thebibliography}

\end{document}
